% !TeX program = xelatex
% !TeX encoding = UTF-8
% Layout specimens only. All text is generic and contains no research claims.
\documentclass[aspectratio=169,11pt]{beamer}
\usetheme{PCL}
\title{中文学术报告模板\\副标题可放在第二行}
\subtitle{版式预览}
\author{汇报人姓名}
\institute{鹏城实验室}
\date{\PCLCurrentYearMonth}

\begin{document}
\PCLTitleFrame
\PCLContentsFrame

\section{单图与双栏版式}
\begin{frame}{单张大图与结论}
  \PCLLead[说明]{上方概括问题，主体区域放置一张清晰的大图。}
  \PCLPlaceholder[40mm]{主图占位\\[2mm]\small 可放入 PDF 矢量图、曲线或系统照片}
  \PCLSource{图注或来源放在此处。此页仅展示版式。}
  \PCLTakeaway{一句话概括本页结论，重点信息可以使用黄色强调。}
\end{frame}

\begin{frame}{图文双栏}
  \begin{columns}[T,onlytextwidth]
    \begin{column}{0.39\textwidth}
      \begin{PCLPanel}{文字说明}
        \begin{itemize}
          \item 第一项核心内容
          \item 第二项方法说明
          \item 第三项结果解释
        \end{itemize}
      \end{PCLPanel}
      \PCLSource{说明文字与右侧图形保持对应。}
    \end{column}
    \begin{column}{0.58\textwidth}
      \PCLPlaceholder[47mm]{图形区域\\[2mm]\small 根据内容调整两栏比例}
    \end{column}
  \end{columns}
  \vfill
  \PCLTakeaway{左侧说明方法，右侧呈现图形或实验结果。}
\end{frame}

\section{对比与公式版式}
\PCLSectionFrame
\begin{frame}{两组结果对比}
  \begin{columns}[T,onlytextwidth]
    \begin{column}{0.485\textwidth}
      \begin{PCLPanel}{条件 A}
        \PCLPlaceholder[32mm]{左图占位}
        \small 图注：简要说明左图。
      \end{PCLPanel}
    \end{column}
    \begin{column}{0.485\textwidth}
      \begin{PCLPanel}{条件 B}
        \PCLPlaceholder[32mm]{右图占位}
        \small 图注：简要说明右图。
      \end{PCLPanel}
    \end{column}
  \end{columns}
  \vfill
  \PCLTakeaway{两图使用一致的尺寸、坐标范围和图例规范。}
\end{frame}

\begin{frame}{公式与图形}
  \PCLLead[模型]{保留关键表达式，并用文字解释符号和条件。}
  \begin{columns}[T,onlytextwidth]
    \begin{column}{0.40\textwidth}
      \begin{PCLPanel}{数学排版示例}
        \[\mathbf{y}=\mathbf{A}\mathbf{x}+\mathbf{n}\]
        \small
        \(\mathbf{x}\)：待估计向量。\par\medskip
        \(\mathbf{A}\)：观测矩阵。\par\medskip
        \(\mathbf{n}\)：噪声向量。
      \end{PCLPanel}
    \end{column}
    \begin{column}{0.57\textwidth}
      \PCLPlaceholder[39mm]{示意图或结果曲线}
    \end{column}
  \end{columns}
  \PCLSource{公式仅用于演示数学排版。}
\end{frame}

\section{总结与表格版式}
\begin{frame}{三项并列总结}
  \begin{columns}[T,onlytextwidth]
    \begin{column}{0.315\textwidth}
      \begin{PCLPanel}[height=43mm]{第一项}
        \textbf{核心观点}\par\medskip
        用两到三行文字说明这一项内容及其意义。
      \end{PCLPanel}
    \end{column}
    \begin{column}{0.315\textwidth}
      \begin{PCLPanel}[height=43mm]{第二项}
        \textbf{核心观点}\par\medskip
        可以替换为一个简洁的示意图或关键数值。
      \end{PCLPanel}
    \end{column}
    \begin{column}{0.315\textwidth}
      \begin{PCLPanel}[height=43mm]{第三项}
        \textbf{核心观点}\par\medskip
        三列的文字层级与视觉重量保持一致。
      \end{PCLPanel}
    \end{column}
  \end{columns}
  \vfill
  \PCLTakeaway{用一句话归纳三项内容之间的联系。}
\end{frame}

\begin{frame}{参数与比较表}
  \PCLLead[参数]{单位、实验条件及统计口径应当标注清楚。}
  \begin{center}
    \renewcommand{\arraystretch}{1.6}
    \begin{tabular}{p{36mm}p{43mm}p{43mm}}
      \toprule
      \textbf{项目} & \textbf{配置 A} & \textbf{配置 B} \\
      \midrule
      参数名称一 & 待填写 & 待填写 \\
      参数名称二 & 待填写 & 待填写 \\
      评价指标 & 待填写 & 待填写 \\
      \bottomrule
    \end{tabular}
  \end{center}
  \PCLSource{注：此处均为模板占位文字，没有实验数值。}
\end{frame}
\PCLClosingFrame
\end{document}
